% ============================================================
%  \section: Speech-Based Conversational Architectures
%  Parent: \section{State of the Art}
% ============================================================

\section{Speech-Based Conversational Architectures}
\label{sec:speech-architectures}

Speech-based conversational agents have to map an incoming acoustic signal
to a spoken response. Ji et al., in the WavChat survey~\cite{ji2024wavechat},
group the architectures used for this task into three families: cascaded,
half-cascade (or semi-cascaded), and end-to-end. The three differ in how
much of the pipeline (speech understanding, language reasoning, speech
synthesis) lives inside a single neural model. The following sections
review each family and discuss the trade-offs.

% ----------------------------------------------------------
\subsection{Cascaded Pipeline (STT--LLM--TTS)}
\label{subsec:cascaded}

The cascaded approach decomposes the problem into three independently
developed stages: an Automatic Speech Recognition (ASR) module transcribes
the user's utterance into text, a Large Language Model (LLM) produces a
textual response, and a Text-to-Speech (TTS) engine renders that response
as audio. Because each component communicates through a well-defined
textual interface, the pipeline is modular: individual stages can be
upgraded, replaced, or scaled without retraining the others.

This architecture underpins the majority of commercially deployed voice
assistants, including the speech pipelines offered by Azure, Google Cloud,
and Amazon Web Services. In the research literature,
X-Talk~\cite{lin2025xtalk} proposes a cascaded dialogue framework targeting
low-latency interaction, while ESPnet-SDS~\cite{li2025espnetsds} provides
an open-source toolkit that implements cascaded spoken-dialogue systems
with pluggable ASR, LLM, and TTS back-ends.

The main advantages of the cascaded pipeline are its maturity and
interpretability. The intermediate text representation makes logging,
debugging, and content moderation straightforward, which is essential in
enterprise and safety-critical settings. Training costs are also low,
since each stage can rely on pre-trained, off-the-shelf models.

The main drawback is latency. The sequential execution of three distinct
inference steps introduces cumulative delays in the range of two to five
seconds, far above the approximately 200\,ms inter-turn gap that
characterises natural human conversation~\cite{heldner2010pauses}.
Streaming reduces this latency in practice. The ASR processes audio while
the user is still speaking, the LLM emits tokens as it generates them, and
the TTS plays each chunk before the full response is ready.
Whisper-Streaming~\cite{machacek2023whisper} reports sub-second incremental
latency on long-form speech using this approach. Errors also compound
across stages: an ASR misrecognition propagates into the LLM prompt and
may trigger an irrelevant response. Finally, the text bottleneck discards
paralinguistic information such as prosody, emotion, and hesitation
markers, which are often crucial for pragmatic interpretation, and prevents
native support for full-duplex interaction.

% ----------------------------------------------------------
\subsection{Half-Cascade (Semi-Cascaded) Architectures}
\label{subsec:half-cascade}

Half-cascade designs fuse two of the three stages and keep the third as
a separate component. The WavChat taxonomy~\cite{ji2024wavechat}
distinguishes three variants, each removing a different inter-stage
boundary.

\emph{Speech-In fusion} (Variant~A) removes the ASR stage: the LLM
ingests audio features directly through a learnable encoder and projection
layer, preserving paralinguistic cues that text transcription would
discard (e.g., SALMONN~\cite{tang2024salmonn},
Freeze-Omni~\cite{wang2024freezeomni}). \emph{Speech-Out fusion}
(Variant~B) absorbs the TTS stage into the LLM, which generates discrete
speech tokens alongside text and so reduces output latency (e.g.,
LLaMA-Omni~\cite{fang2024llamaomni}). \emph{Inner Monologue} (Variant~C)
keeps text as an internal reasoning scaffold while producing speech output
directly, combining coherent reasoning with expressive synthesis (e.g.,
SpeechGPT~\cite{zhang2023speechgpt}).

These architectures cut latency and use richer multimodal representations
than cascaded pipelines, but at a cost. Aligning audio and text
representation spaces needs carefully staged training, and the model's
language abilities can degrade after fine-tuning on speech data, an effect
known as reasoning degradation~\cite{chen2025urobench}. Speech quality
also tends to fall behind dedicated TTS systems, and none of these models
has been evaluated beyond dyadic interaction.

% ----------------------------------------------------------
\subsection{End-to-End Speech-to-Speech Models}
\label{subsec:e2e}

End-to-end architectures collapse the entire spoken-dialogue pipeline into
a single neural model that maps audio input directly to audio output. A
key enabler is the family of \emph{neural audio codecs} (SoundStream,
EnCodec, Mimi), which convert continuous waveforms into compact sequences
of discrete tokens suitable for autoregressive modelling.

Moshi~\cite{defossez2024moshi} was the first publicly demonstrated
full-duplex speech-to-speech model, with a mouth-to-ear latency of about
200\,ms. Commercial systems such as GPT-4o (OpenAI), Gemini~2.5 (Google),
Qwen3-Omni (Alibaba), and Kimi Audio (Moonshot AI) have since adopted
end-to-end or near-end-to-end designs, although their architectural
details remain mostly proprietary. In the academic domain, dGSLM extends
the Generative Spoken Language Model to dialogue by jointly modelling two
audio channels.

End-to-end models have a clear appeal: removing the intermediate
representations cuts latency and preserves paralinguistic cues (pitch
contour, speaking rate, emotional colouring) that cascaded systems
discard. They also support full-duplex interaction natively, since the
model can listen and speak at the same time.

Current end-to-end systems still have important limitations.
URO-Bench~\cite{chen2025urobench} shows that, on reasoning-intensive
tasks, end-to-end speech models perform substantially worse than their
text-only LLM counterparts, suggesting that the speech modality introduces
a reasoning bottleneck. Training costs are high, since they need large
volumes of paired speech-to-speech data. Reliability is also a concern:
Lin et al.~\cite{lin2025xtalk} observe that end-to-end speech models are
``often unsuitable for deployment'' in production settings due to
hallucination, inconsistent output quality, and limited controllability.

% ----------------------------------------------------------
% Closing paragraph: ties the architecture analysis to the multi-party
% setting that motivates AIutami's cascaded choice. Not a separate
% subsection on purpose, to avoid the "research gap" framing — the
% chapter-level gap discussion lives in the final Discussion section.

Every multi-party speech-based AI system reported in the literature is
built on a cascaded pipeline. The ARI robot~\cite{addlesee2024ari} and
the Furhat platform~\cite{axelsson2025furhat} both use cascaded STT, LLM,
and TTS to handle speaker identification, overlapping speech, and group
turn-taking. End-to-end and half-cascade models have so far only been
tested in dyadic settings. This is an empirical observation, not a
theoretical limit: nothing prevents a neural speech-to-speech model from
supporting groups in principle, but no published work shows it actually
working. A multi-party system built today therefore commits to the
cascaded pattern, and the interesting design choices sit at the layers
above the pipeline: turn-taking, addressee detection, and moderation.
These are the focus of the rest of this chapter.
